\documentclass[11pt,a4paper]{article}
\usepackage{amsmath,amssymb,amsthm}
\usepackage[margin=2.5cm]{geometry}
\usepackage{hyperref}

\newtheorem{definition}{Definition}[section]
\newtheorem{theorem}[definition]{Theorem}
\newtheorem{proposition}[definition]{Proposition}
\newtheorem{remark}[definition]{Remark}
\newtheorem{conjecture}[definition]{Conjecture}

\title{Hyperseed Formalization:\\
  Parafinity}
\author{ProtomegaTron (automated formalization)}
\date{2026-09-18}

\begin{document}
\maketitle

\begin{abstract}
We formalize Ben Goertzel's ``Parafinity'' (2022-01-15) within the
Hyperseed ontology. Key mappings: parafinity as the fiber boundary zone,
constructive nonstandard analysis as computable fiber, and AGI reasoning
as parafinite fiber operations.
\end{abstract}

\tableofcontents

\section{The parafinite zone: fiber boundary}

\begin{theorem}[Parafinity = fiber boundary --- claims 4--9, 13]
The parafinite zone is where fiber meets base --- too structured for raw
base treatment (finite enumeration) but not fully abstract (classical
infinity):
\[
  \text{finite} \subset \text{parafinite} \subset \text{infinite}
  \quad \text{with each } \subsetneq
\]
Parafinite structures have qualitatively different properties from both
ordinary finite and classical infinite --- they constitute a genuine
third zone, not just ``large finite'' or ``small infinite.''

In Hyperseed terms, fiber operations naturally live in the parafinite
zone: the fiber is too complex for finite enumeration of all fibers over
all base points, but structured enough (via the d-calculus) for
constructive manipulation.
\end{theorem}

\section{Computable fiber: constructive approach}

\begin{proposition}[Constructive = computable --- claims 1--3, 14]
Mycielski's constructive nonstandard analysis provides computable
infinitesimals --- infinitesimal-like objects that can actually be
computed with:
\[
  \epsilon_{\text{constructive}} \approx 0 \quad \text{but} \quad
  \text{computable}(\epsilon_{\text{constructive}}) = \text{true}
\]
This is essential for Hyperseed implementation: the fiber must be
constructively manipulable, not just abstractly defined. Computable
fiber is fiber with effective operations --- actual algorithms for
fiber transport, composition, and evolution.
\end{proposition}

\section{AGI in the parafinite zone}

\begin{conjecture}[AGI reasoning is parafinite --- claims 10--12, 15--16]
AGI reasoning naturally operates in the parafinite zone:
\[
  \text{AGI} \in \text{Parafinite}(\text{Fiber})
  = \{F : \text{finite} \subsetneq F \subsetneq \text{infinite}\}
\]
The d-calculus may be naturally parafinite --- operating on structures
too complex for finite enumeration but constructively manipulable.
Paraconsistency is especially natural in this zone, where classical
distinctions (finite/infinite, consistent/inconsistent) soften into
gradients.
\end{conjecture}

\section{Cross-references}

\begin{remark}[Connections]
\begin{itemize}
\item \textbf{Paraconsistent Interzones} (2021-08-13): paraconsistency
  as interzone --- boundary zone where classical distinctions soften.
\item \textbf{The Mind-Blowing Paraconsistency} (2021-09-23): deep dive
  into paraconsistent logic.
\item \textbf{Tensor Logic} (2025-12-16): tensor operations as
  computable fiber --- same constructive imperative.
\item \textbf{Facing the Meta-Abstracted Dragon} (2022-07-31): infinite
  tower of meta-abstraction --- a parafinite structure.
\end{itemize}
\end{remark}

\end{document}
